\section{implementatie}
\label{sec:implementatie}

Dit hoofdstuk beschrijft de implementatie van het sensor fusion systeem. De realisatie is stapsgewijs opgebouwd: de configuratie van de ontwikkelomgeving, de IMU integratie inclusief datacommunicatie, en tot slot de realisatie van het Extended Kalman Filter (EKF).

\subsection{Ontwikkelomgeving en testopstelling}

Voor de eindgebruiker is het systeem ontworpen om draadloos firmware te ontvangen via Over The Air (OTA) updates. Aangezien dit niet efficiënt is voor programmeren, is een fysieke testopstelling gerealiseerd met behulp van een SEGGER J-Link debug probe \cite{segger_rtt}. \\

Hoewel de J-Link standaard een 20 pins JTAG-connector heeft, communiceert de microcontroller op de printplaat via de Serial Wire Debug (SWD) interface. Omdat de printplaat niet is voorzien van een fysieke debug-header, is de verbinding handmatig tot stand gebracht. De PCB bevat specifieke testpads bedoeld voor de vier SWD signalen (VCC, GND, SWDIO en SWCLK) waaraan tijdelijk 4 kabels zijn gesoldeerd. Deze staan direct in verbinding met de J-Link (zie Figuur \ref{fig:JTAG}).

\begin{figure}[htbp]
    \centering
    \includegraphics[width=0.2\linewidth]{JTAG_Pinout.png}
    \caption{Pinout van de JTAG connector waarbij pinnen 1, 7, 9 en GND zijn gebruikt voor de SWD-verbinding \cite{microchip-tseg-jlink}.}
    \label{fig:JTAG}
\end{figure}

\subsection{IMU integratie}

Voor de ASM330LHH IMU moeten diverse parameters worden geconfigureerd om te voldoen aan de dynamische eisen van het systeem. De functionele keuzes zijn samengevat in Tabel~\ref{tab:imu_config}. De vertaalslag hiervan naar de exacte bit-waarden in de controleregisters (zoals \texttt{CTRL1\_XL}) is gedocumenteerd in Bijlage X.

\begin{table}[htbp]
    \centering
    \caption{Geconfigureerde functionele parameters van de ASM330LHH IMU.}
    \label{tab:imu_config}
    \renewcommand{\arraystretch}{1.3}
    \begin{tabularx}{\textwidth}{|X|X|}
        \hline
        \rowcolor{green!20}
        \textbf{Parameter} & \textbf{Gekozen waarde} \\ \hline
        Accelerometerbereik & $\pm$4 g \\ \hline
        Gyroscoopbereik & $\pm$500 DPS \\ \hline
        Outputfrequentie (ODR) & 208 Hz \\ \hline
        Laagdoorlaatfilter & 50 Hz \\ \hline
        Block Data Update (BDU) & Ingeschakeld \\ \hline
    \end{tabularx}
\end{table}

Om de beweging accuraat vast te leggen zonder clipping, is het accelerometerbereik ingesteld op $\pm$4 g. Dit geeft een veiligheidsmarge van 60\% ten opzichte van de berekende piekversnelling van 2,5 g \cite{stmicro_asm330lhh}. Voor het gyroscoopbereik is gekozen voor $\pm$500 DPS \cite{stmicro_asm330lhh}. Ten opzichte van de eerder berekende worst-case hoeksnelheid van 37 DPS op de renbaan, biedt dit een ruime marge (factor 14) voor abrupte en onverwachte bewegingen. \\

Voor de outputfrequentie (ODR) is een waarde van 208 Hz geconfigureerd. Dit sluit aan bij de standaard die gehanteerd wordt in het EquiMoves systeem voor paarden \cite{InertiaEquineGait} en het onderzoek naar IMU oriëntatieschatting bij hardlopers \cite{Fan2025Influence}. De IMU beschikt bovendien over een digitale laagdoorlaatfilter die in het systeem is afgesteld op 50 Hz. Hiervoor is gekozen, omdat onderzoek concludeert dat bewegingsfrequenties bij renpaarden zich onder de 10 Hz bevinden \cite{witte2006effect}. Hiermee wordt hoogfrequente ruis gedempt, zonder het werkelijke signaal aan te tasten. \\

Als laatst is de \textit{Block Data Update} (BDU) ingeschakeld om fouten bij het uitlezen van de seriële interface te voorkomen. Zoals geadviseerd door STMicroelectronics \cite{ST_AN5296} (pagina 18), garandeert dit dat de hoge en lage bytes van de outputregisters afkomstig zijn van hetzelfde meetsample.


\subsubsection{Gyroscooop kalibratie}
Wanneer de IMU goed is geconfigureerd kan gekeken worden naar het uitlezen van de data. Voor het uitlezen van de data van de accelerometer en de gyroscoop moet gekeken worden naar registers OUTX\_L\_A en OUTX\_L\_G. Deze registers bevatten meerdere 16 bit words met de waardes van de x, y en z as voor de accelerometer en de pitch, roll en yaw voor de gyroscoop. Waardes die uit deze registers verkregen zijn, zijn te zien in figuur \ref{fig:sensorwaardes}.

\begin{figure}[h]
    \centering
    \includegraphics[width=1\textwidth]{sensorwaardes_origineel.png}
    \caption{Gyroscoop en accelerometer waardes bij stilstand op een tafel. Hierbij laat het bovenste venster de gyroscoop draaisnelheden zien en het onderste venster de g-krachten die de accelerometer waarneemt. De zwarte lijn is bedoeld om aan te tonen wat de totale g-kracht is die zich op het apparaat bevindt}
    \label{fig:sensorwaardes}
\end{figure}

Hierin zijn de waardes te zien van het apparaat die stil ligt op een tafel. Hierin zijn de verschillende g-krachten op de assen te zien en die resulteren gezamenlijk tot ongeveer 1000 milli g wat overeenkomt met de zwaartekracht. Er is duidelijk een offset te zien en de waardes schommelen. Het valt ruim binnen de specificaties van de IMU aangezien de datasheet vermeld dat het bij een zero rate level tussen de -10 en +10 dps kan zitten. De accelerometer is daarentegen vrij gelijk. Om met de waardes van de gyroscope te werken moet ernaar gekeken worden om de waardes richting 0 te krijgen en de schommelingen te verminderen. Om de schommelingen te verminderen is een filter toegevoegd die zich binnen de configuratie van de IMU bevindt. De resultaten hiervan zijn te zien in figuur \ref{fig:sensorwaardes_filter}.

\newpage

\begin{figure}[h]
    \centering
    \includegraphics[width=1\textwidth]{sensorwaardes_filter.png}
    \caption{De filter configuratie van de IMU is gebruikt voor de gyroscoop. De schommelingen zijn hierdoor aanzienlijk minder}
    \label{fig:sensorwaardes_filter}
\end{figure}

Zoals te zien in de afbeelding, hebben MEMS gyroscopen (Micro-Electro-Mechanical System) een systematische meetafwijking bij stilstand, ook wel \textit{zero-rate offset} of bias genoemd. Deze offset zal accumuleren bij het bepalen van de richting als dit niet gecompenseerd wordt. Om dit te verhelpen wordt bij het opstarten een eenmalige kalibratie uitgevoerd. De eerste 400 samples worden niet gebruikt vanwege de pieken van de sensor. Vervolgens wordt een gemiddelde gebruikt van 200 opeenvolgende samples als kalibratieoffset. De berekende offset wordt bij iedere gyroscoopuitlezing afgetrokken van de ruwe sensorwaarde. Het resultaat hiervan is te zien in figuur \ref{fig:sensorwaardes_kalibratie}.

\begin{figure}[h]
    \centering
    \includegraphics[width=1\textwidth]{sensorwaardes_kalibraties.png}
    \caption{De offset van de gyroscoop waardes zijn van de metingen afgetrokken}
    \label{fig:sensorwaardes_kalibratie}
\end{figure}

\subsubsection{SPI-busprobleem bij de IMU-uitlezing}
 
Bij langere testmomenten bleek dat de sensordata sporadisch incorrect was. De uitgelezen waardes bevatten uitschieters die niet overeenkwamen met de werkelijke situatie. Dit is duidelijk te zien in figuur \ref{fig:imu_pieken}. De pieken waren consistent en werden vrijwel altijd veroorzaakt op momenten dat andere apparaten ook gebruik maakten van de SPI-bus. Op de tracker delen de IMU, de seriële flash en de CAN controller dezelfde SPI-bus. Dit betekent dat alle drie de apparaten via dezelfde datalijn (MOSI/MISO) en kloklijnen communiceren met de microcontroller. \\

\begin{figure}[h]
    \centering
    \includegraphics[width=1\textwidth]{imu_pieken.png}
    \caption{Willekeurige grote pieken in de gyroscoop data. Dit zijn resultaten op een stilstaande tafel}
    \label{fig:imu_pieken}
\end{figure}

Het probleem ontstaat doordat elk apparaat op de SPI-bus een eigen Chip Select (CS) pin heeft waarmee de microcontroller aangeeft met welk apparaat hij wil communiceren. Als de CS-pin van de IMU laag is, luistert de IMU naar alle data op de bus en hiermee dus ook data die bestemd is voor de flash of de CAN-controller. Wanneer de CS-pin van de IMU niet correct hoog werd gezet na een transactie, of wanneer een andere transactie startte terwijl de IMU nog actief was, kon de IMU incorrecte data vertonen. \\

De oplossing zit in het beheren van de CS pinnen bij elke transactie. Per apparaat wordt de pin direct voor de dataoverdracht laag gezet en direct erna weer hoog. 

\begin{lstlisting}[language=C]
gpio_pin_clear(config->cs_pin);   // CS laag: selecteer IMU
spi_xfer(config->spi, &xfer);     // voer transactie uit
gpio_pin_set(config->cs_pin);     // CS hoog: deselecteer IMU
\end{lstlisting}

Bovendien verwacht de ASM330LHH een soort dummy overdracht voordat de daadwerkelijke communicatie plaatsvindt. Dit heeft te maken met het SPI protocol. Bij SPI vindt dataoverdracht gelijktijdig plaats. Zodra de microcontroller (master) de SCK activeert om een commando te versturen, wordt er op hetzelfde moment een byte vanuit de sensor (slave) teruggestuurd. De sensor heeft nog niet de tijd gehad om de binnenkomende commando te verwerken. Hierdoor ontstaat onbruikbare data. Door bewust eerst een dummy-byte van \texttt{0xFF} te sturen zonder de ontvangen byte te verwerken, wordt deze initiële toestand weggegooid. Hierna bevinden zowel de sensor als de microcontroller zich in een bekende begintoestand. Na het toepassen van deze methodes zijn de sporadische uitschieters verdwenen uit de sensordata.


\subsection{Extended Kalman Filter}





\begin{comment}
    
Het EKF combineert de IMU-data met de RTK data tot één betrouwbare positieschatting. Dit filter is opgebouwd uit meerdere onderdelen waarvan deze zijn weergegeven in tabel \ref{tab:ekf_componenten}.

\newpage

\begin{table}[htbp]
    \centering
    \caption{Overzicht van de componenten van het Extended Kalman Filter.}
    \label{tab:ekf_componenten}
    \renewcommand{\arraystretch}{1.3}
    \begin{tabularx}{\textwidth}{|l|l|X|}
        \hline
        \rowcolor{green!20}
        \textbf{Component} & \textbf{Frequentie} & \textbf{Functie} \\ \hline
        Complementary filter & 200 Hz & 
        Bepaalt pitch en roll uit gyroscoop en accelerometer voor 
        zwaartekrachtcompensatie \\ \hline
        State vector $\mathbf{x}$ & n.v.t & 
        Bevat alle geschatte grootheden: positie, snelheid, heading en 
        sensorbiassen \\ \hline
        Motion model + Jacobiaan $F$ & 200 Hz & 
        Voorspelt de nieuwe toestand op basis van IMU-data 
        (Updatestap) \\ \hline
        Procesruismatrix $Q$ & 200 Hz & 
        Modelleert de onzekerheid van het bewegingsmodel per tijdstap \\ \hline
        Covariantiematrix $P$ & 200 Hz & 
        Houdt bij hoe zeker het filter is over elke grootheid \\ \hline
        Meetmodel $H$ + R-matrix & 25 Hz & 
        Beschrijft welke grootheden worden gemeten en hoe onzeker die 
        meting is \\ \hline
        Kalman gain $K$ & 25 Hz & 
        Bepaalt hoeveel gewicht de GNSS-meting krijgt ten opzichte van 
        de updatestap \\ \hline
        NIS-check & 25 Hz & 
        Toetst statistische plausibiliteit van de GNSS-meting; verwerpt 
        uitschieters \\ \hline
        ZUPT & Continu & 
        Corrigeert snelheidsschatting naar nul bij gedetecteerde 
        stilstand \\ \hline
    \end{tabularx}
\end{table}

De updatestap wordt elke 5\,ms uitgevoerd op basis van IMU-data. De updatestap 
wordt uitgevoerd zodra een nieuwe GNSS-meting beschikbaar is (elke 40\,ms). De 
onderdelen worden hieronder stapsgewijs toegelicht.


\subsubsection{Complementary filter voor attitude}

Voordat de IMU-data het EKF binnenkomt, wordt de ruwe accelerometerdata en gyroscoopdata eerst verwerkt door een complementary filter. Dit filter bepaalt de hellingshoek van de tracker, wat nodig is om de juiste rijversnelling uit de sensordata te halen. \\

Een accelerometer meet namelijk niet alleen de rijversnelling, maar ook de zwaartekracht. Een accelerometer die stilstaat op een tafel meet een versnelling van 1 g, ook wel de zwaartekracht. Zolang de tracker perfect horizontaal ligt, valt die zwaartekracht volledig op de Z-as van de sensor en heeft dit geen invloed op de X-as. Maar als de tracker licht gekanteld is gemonteerd, valt de zwaartekracht niet meer puur op de Z-as. Een deel ervan lekt dan door naar de X-as, waardoor de sensor een constante versnelling op die as meet terwijl het voertuig stilstaat. Het filter zou dit onjuist als rijversnelling interpreteren. \\

Om dit te corrigeren wordt ten eerste bepaald hoe de tracker gekanteld is. Een sensor heeft drie assen en kan om elk van die assen draaien. De kanteling naar voren of achteren wordt de \textit{pitch} genoemd en de kanteling naar links of rechts de \textit{roll}. Samen beschrijven pitch en roll hoe de tracker in de ruimte staat ten opzichte van het horizontale vlak. \\

Het complementary filter bepaalt deze hoeken door de gyroscoop en de accelerometer te combineren. De gyroscoop meet hoeksnelheden en kan door die te integreren over de tijd bijhouden hoeveel de tracker is gekanteld. Dit werkt goed voor snelle bewegingen, maar door kleine meetfouten die zich opstapelen raakt de schatting over langere tijd langzaam uit koers. De accelerometer heeft dit probleem niet aangezien deze bij stilstand altijd precies in de richting wijst van de zwaartekracht. Dit biedt een absolute referentie voor de hellingshoek. Het nadeel is dat de accelerometer gevoelig is voor trillingen tijdens het rijden, waardoor die referentie bij beweging minder betrouwbaar is. Door beide sensoren samen te gebruiken worden hun zwaktes gecompenseerd. De gyroscoop zorgt voor een snelle en vloeiende respons en de accelerometer corrigeert de langzame drift. In formule \ref{eq:complementary} is dit uiteengezet.

\begin{equation}
    \theta_{\text{pitch}} = 0{,}98 \cdot \left(\theta_{\text{pitch, prev}} + 
    \omega_{\text{pitch}} \cdot \Delta t\right) + 0{,}02 \cdot 
    \arctan\!\left(\frac{-a_x}{\sqrt{a_y^2 + a_z^2}}\right) \cdot \frac{180}{\pi}
    \label{eq:complementary}
\end{equation}

Hierbij is $\omega_{\text{pitch}}$ de gemeten hoeksnelheid, zijn $a_x$, $a_y$ en $a_z$ de gemeten versnellingen op de drie assen en is $\Delta t$ de tijdstap. De coëfficiënt 0{,}98 bepaalt hoeveel gewicht de gyroscoop krijgt ten opzichte van de accelerometer. Hetzelfde wordt ook berekend voor de rollhoek. \\


\subsubsection{Zwaartekrachtcompensatie}

Zodra de hellingshoek (pitch) bekend is vanuit het complementary filter, wordt 
deze gebruikt om de zwaartekrachtcomponent weg te rekenen uit de gemeten 
versnelling op de x-as. Wanneer een sensor gekanteld is, projecteert de 
zwaartekracht zich gedeeltelijk op de meetassen van de accelerometer. Zoals 
beschreven in de applicatienote van Analog Devices~\cite{analogdevices_an1057}, 
geldt dat de projectie van de zwaartekracht op de x as gelijk is aan 
$g \cdot \sin(\theta_p)$. Bij een pitch
$\theta_p$ projecteert de zwaartekracht zich dus als $-\sin(\theta_p)$ op 
de x as van de sensor. Dit component wordt afgetrokken van de ruwe meting 
zoals uitgedrukt in formule~\ref{eq:grav_comp}.

\begin{equation}
    a_{\text{net}} = a_x - (-\sin\theta_p) = a_x + \sin\theta_p
    \label{eq:grav_comp}
\end{equation}

Hierbij is $a_x$ de ruwe accelerometerwaarde in g en $\theta_p$ de 
pitchhoek in radialen. De netto versnelling wordt vervolgens omgezet 
naar m/s² in formule~\ref{eq:accel_ms2}.

\begin{equation}
    a_{\text{ms}^2} = a_{\text{net}} \cdot 9{,}81
    \label{eq:accel_ms2}
\end{equation}

In code is dit als volgt geïmplementeerd:

\begin{lstlisting}[language=C, label=lst:gravity_comp]
// Omzetten van milli-g naar g
float raw_ax = (float)current_accel.x / 1000.0f;

// Pitchhoek omzetten naar radialen
float pitch_rad = ((float)current_attitude.pitch_x100 / 100.0f)
                  * ((float)M_PI / 180.0f);

// Zwaartekrachtcomponent op de X-as bepalen en aftrekken
float grav_x    = -sinf(pitch_rad);
float net_ax    = raw_ax - grav_x;

// Omzetten van g naar m/s2
float accel_ms2 = net_ax * 9.81f;
\end{lstlisting}

Deze methode gebruikt alleen de pitch omdat de tracker zo is 
gemonteerd dat de rijrichting overeenkomt met de x as van de sensor. 
De roll heeft in dat geval geen noemenswaardige invloed op de 
voorwaartse versnelling en wordt daarom buiten beschouwing gelaten. 
Dit is een vereenvoudiging van de hiervoor uitgewerkte benadering 
waarbij door middel van rotatiematrices de versnellingsvector volledig 
naar het horizontale vlak wordt teruggedraaid~\cite{nxp2013}. Deze 
methode is opgenomen in bijlage~\ref{app:zwaartekracht} voor de 
volledigheid.

\subsubsection{State vector}

Om het filter op te bouwen moeten de grootheden van het systeem bijgehouden worden. Dit wordt bereikt door middel van de state vector. Dit is een lijst van alle variabelen die het filter op elk tijdstip schat. In dit geval bestaat de state vector uit zes elementen zoals te zien in formule \ref{eq:toestandsvector}

\begin{equation}
    \mathbf{x} = \begin{bmatrix} X & Y & v & \theta & b_a & b_g \end{bmatrix}^T
    \label{eq:toestandsvector}
\end{equation}

Hierbij zijn $X$ en $Y$ de positie in meters in het lokale coördinatenstelsel, $v$ de rijsnelheid in m/s en $\theta$ de heading (rijrichting) in radialen. Deze vier grootheden vormen de kern van de positieschatting en worden direct gebruikt voor de routebepaling. \\

Naast deze vier grootheden bevat de state vector twee biastermen: $b_a$ voor de accelerometer en $b_g$ voor de gyroscoop. Een bias is een kleine systematische meetafwijking die aanwezig is onafhankelijk van de werkelijke beweging. Hoewel bij het opstarten al een eenmalige gyroscoop kalibratie wordt uitgevoerd die de zero-rate offset verwijdert, blijft er altijd een kleine resterende fout over. Deze fout kan langzaam veranderen tijdens gebruik door drift of eventueel de opwarming van de sensor. \\

Door $b_a$ en $b_g$ op te nemen in de state vector schat het EKF deze resterende afwijking continu bij tijdens het rijden en trekt deze af van de gemeten versnelling en hoeksnelheid voordat ze worden geïntegreerd. Dit wordt in opvolgende paragrafen duidelijker toegelicht. Zonder deze correcties zou een kleine constante fout in de gemeten versnelling zich bij integratie opstapelen tot een toenemende fout in snelheid. \\

Het filter wordt geïnitialiseerd zodra de eerste RTK fix plaatsvindt, omdat dit het moment is dat het systeem het meest nauwkeurig is. Het EKF zal dit punt gebruiken als positieoorsprong van het te gebruiken coördinatenstelsel.
 
\subsubsection{Motion model}

Na het definiëren van de state vector, moet het filter weten hoe zijn toestand (state) verandert over de tijd. Dit wordt ook wel het motion model genoemd. Dit model beschrijft welke waarde elke grootheid naar verwachting heeft na elk tijdstip, op basis van de IMU metingen. \\

Om dit te bereiken wordt een zogenoemde \textit{unicycle model} gebruikt zoals beschreven in paragraaf 13.1.2.3 van \textit{`Planning Algorithms'} \cite{lavalle2006planning}. Dit is een standaard rijmodel binnen de robotica en besturingstechniek dat ervan uitgaat dat het object beweegt met snelheid $v$ in de richting $\theta$ en draait met een hoeksnelheid $\dot{\theta}$ (yaw-rate). Zijwaartse slip wordt in dit geval niet volledig meegenomen.
 
\begin{align}
    X_{k+1}      &= X_k + v_k \cdot \cos(\theta_k) \cdot \Delta t \label{eq:fx} \\
    Y_{k+1}      &= Y_k + v_k \cdot \sin(\theta_k) \cdot \Delta t \label{eq:fy} \\
    v_{k+1}      &= v_k + (a_{\text{meas}} - b_{a,k}) \cdot \Delta t \label{eq:fv} \\
    \theta_{k+1} &= \theta_k + (\omega_{\text{meas}} - b_{g,k}) \cdot \Delta t \label{eq:ftheta} \\
    b_{a,k+1}    &= b_{a,k} \label{eq:fba} \\
    b_{g,k+1}    &= b_{g,k} \label{eq:fbg}
\end{align}
 
De eerste twee vergelijkingen (\ref{eq:fx} en \ref{eq:fy}) beschrijven hoe de oost ($X_{k+1}$) en noord ($Y_{k+1}$) positie verandert op basis van de huidige snelheid en rijrichting. Vergelijking \ref{eq:fv} werkt de snelheid bij op basis van de gemeten versnelling, waarbij de accelerometerbias als correctie er af gehaald wordt. Vergelijking \ref{eq:ftheta} werkt de heading bij op basis van de gemeten yaw rate (zie figuur \ref{fig:oriëntatie}), die vervolgens ook gecorrigeerd wordt door de gyroscoop bias. \\

De laatste twee vergelijkingen (\ref{eq:fba} en \ref{eq:fbg}) tonen de biassen die over één tijdstap van 5\,ms worden bepaalt. Deze blijft gedurende deze tijdstap constant, omdat deze termen afhankelijk zijn van de updatestap. \\

\subsubsection{Covariantiematrix en Jacobian matrix}

Naast de toestand zelf houdt het EKF bij hoe zeker het is over elke 
schatting. Dit wordt bijgehouden in de covariantiematrix $P$, een 
$6 \times 6$ matrix~\cite{welch_kalman}. Een concreet voorbeeld: als 
het filter net is geïnitialiseerd en nog weinig GNSS-correcties heeft 
ontvangen, is de onzekerheid over de positie groot. Dit wordt 
weerspiegeld door grote waarden op de diagonaal van $P$ voor de 
positie-elementen $P[0][0]$ en $P[1][1]$. Naarmate meer 
GNSS-correcties binnenkomen en het filter convergeert, nemen deze 
waarden af. Een groot $P[0][0]$ betekent dus concreet: het filter 
heeft weinig vertrouwen in zijn huidige schatting van de 
oost-positie $X$. \\

De off-diagonaalelementen van $P$ beschrijven de correlatie tussen 
twee grootheden. Een niet-nul waarde op positie $P[0][4]$ betekent 
dat de onzekerheid in de oost-positie $X$ en de onzekerheid in de 
accelerometerbias $b_a$ aan elkaar gekoppeld zijn. Dit is een 
belangrijk mechanisme dat verderop wordt uitgelegd. \\

Omdat vergelijkingen~\ref{eq:fx} en~\ref{eq:fy} niet-lineair zijn 
in $\theta$ — ze bevatten een sinus en cosinus — kan $P$ niet 
rechtstreeks worden bijgewerkt. Het EKF lost dit op door het 
bewegingsmodel bij elke tijdstap lokaal te lineariseren via de 
Jacobiaan $F$~\cite{thrun2005probabilistic, bar_shalom_estimation}. 
De Jacobiaan is de matrix van alle partiële afgeleiden van het 
bewegingsmodel naar de toestandsgrootheden. Een partiële afgeleide 
geeft aan hoeveel een uitkomst verandert als één invoergrootheid een 
kleine stap neemt terwijl de rest constant blijft. \\

Als concreet voorbeeld: vergelijking~\ref{eq:fx} beschrijft dat 
$X_{k+1} = X_k + v_k \cdot \cos(\theta_k) \cdot \Delta t$. De 
partiële afgeleide naar $v$ geeft aan hoeveel $X_{k+1}$ verandert 
als $v$ een kleine stap neemt. Dat is $\cos(\theta) \cdot \Delta t$ 
— precies de term die op rij 0, kolom 2 van de Jacobiaan staat. De 
partiële afgeleide van diezelfde vergelijking naar $\theta$ is 
$-v \cdot \sin(\theta) \cdot \Delta t$, wat op rij 0, kolom 3 staat. 
Op dezelfde manier worden alle andere afgeleiden bepaald voor de 
overige bewegingsvergelijkingen~\ref{eq:fy}--\ref{eq:fbg}. Dit 
levert de volgende Jacobiaan op die elke tijdstap opnieuw wordt 
berekend op basis van de huidige schatting van $v$ en $\theta$:

\begin{equation}
F = \begin{bmatrix}
1 & 0 & \cos\theta \cdot \Delta t & -v \sin\theta \cdot \Delta t & 0 & 0 \\
0 & 1 & \sin\theta \cdot \Delta t &  v \cos\theta \cdot \Delta t & 0 & 0 \\
0 & 0 & 1 & 0 & -\Delta t & 0 \\
0 & 0 & 0 & 1 & 0 & -\Delta t \\
0 & 0 & 0 & 0 & 1 & 0 \\
0 & 0 & 0 & 0 & 0 & 1
\end{bmatrix}
\label{eq:jacobiaan_F}
\end{equation}

De structuur van de matrix vertelt direct hoe de grootheden met 
elkaar verbonden zijn. De enen op de diagonaal geven aan dat elke 
grootheid zichzelf doorgeeft naar de volgende tijdstap. De nulwaarden 
geven aan dat die twee grootheden elkaar niet direct beïnvloeden — 
zo heeft de oost-positie $X$ geen invloed op de heading $\theta$ en 
andersom. 

Twee elementen verdienen extra aandacht. Op rij 2 (snelheid $v$) 
staat $-\Delta t$ op kolom 4 (accelerometerbias $b_a$). Dit zegt: 
een kleine toename van de bias $b_a$ leidt na één tijdstap tot een 
afname van de snelheidsschatting met $\Delta t$, precies zoals 
vergelijking~\ref{eq:fv} beschrijft. Op rij 3 (heading $\theta$) 
staat eveneens $-\Delta t$ op kolom 5 (gyroscoop-bias $b_g$), 
overeenkomstig vergelijking~\ref{eq:ftheta}.

Wanneer $F \cdot P \cdot F^T$ wordt berekend verspreidt deze 
koppeling zich door $P$. Na een aantal tijdstappen zonder 
GNSS-correctie is er een niet-nul waarde opgebouwd op 
$P[0][4]$ — de correlatie tussen de oost-positie en de 
accelerometerbias. Dit betekent dat het filter weet dat een deel 
van de positieonzekerheid wordt veroorzaakt door de bias. Zodra 
er een GNSS-correctie binnenkomt die een positiefout signaleert, 
gebruikt het filter via de Kalman gain deze correlatie om een deel 
van die positiefout toe te schrijven aan de bias en die dienovereenkomstig 
bij te stellen. Zo werkt de biascorrectie indirect via de 
covariantiematrix.




















\begin{comment}

Aangezien het systeem niet lineair is, zoals te zien is aan de sinus en cosinus in de vergelijkingen, kan $P$ niet rechtstreeks worden bijgewerkt. Het EKF lost 
dit op door het bewegingsmodel bij elke tijdstap lokaal te lineariseren via 
de Jacobian matrix $F$. Dit is een matrix die beschrijft hoe een kleine verandering 
in één grootheid doorwerkt op alle andere grootheden na één tijdstap. De 
Jacobian matrix wordt elke tijdstap opnieuw berekend op basis van de huidige 
schatting. Deze matrix is te zien in figuur \ref{eq:jacobiaan_F}.

\begin{equation}
F = \begin{bmatrix}
1 & 0 & \cos\theta \cdot \Delta t & -v \sin\theta \cdot \Delta t & 0 & 0 \\
0 & 1 & \sin\theta \cdot \Delta t &  v \cos\theta \cdot \Delta t & 0 & 0 \\
0 & 0 & 1 & 0 & -\Delta t & 0 \\
0 & 0 & 0 & 1 & 0 & -\Delta t \\
0 & 0 & 0 & 0 & 1 & 0 \\
0 & 0 & 0 & 0 & 0 & 1
\end{bmatrix}
\label{eq:jacobiaan_F}
\end{equation}

Met de Jacobiaan wordt $P$ bijgewerkt in de updatestap. De term 
$F_k \cdot P_k \cdot F_k^T$ beschrijft hoe de huidige onzekerheid 
doorwerkt na één tijdstap. Omdat het bewegingsmodel nooit perfect is, 
wordt daarbovenop de procesruismatrix $Q$ opgeteld. Dit modelleert dat 
de werkelijkheid altijd iets afwijkt van wat de vergelijkingen voorspellen:

\begin{equation}
    P_{k+1}^- = F_k \cdot P_k \cdot F_k^T + Q
    \label{eq:P_predict}
\end{equation}

Door afrondingsfouten in zwevendekommaberekeningen kan $P$ na verloop van 
tijd licht asymmetrisch worden, wat de stabiliteit van het filter kan 
verstoren. Om dit te voorkomen wordt $P$ na elke stap gesymmetriseerd:

\begin{equation}
    P \leftarrow \frac{P + P^T}{2}
    \label{eq:symmetrisatie}
\end{equation}

Tot slot wordt na elke updatestap de heading teruggevouwen naar het 
bereik $[-\pi, +\pi]$. Wanneer het systeem meerdere rondjes rijdt, kan de 
geïntegreerde heading onbeperkt oplopen — bijvoorbeeld naar $7\pi$. Als 
het filter dan een GNSS-meting van $\pi$ ontvangt, berekent het een grote 
correctie terwijl de werkelijke fout nul is. Door de heading na elke stap 
terug te vouwen wordt dit voorkomen:

\begin{lstlisting}[language=C]
while (s_status[3][0] >  M_PI) s_status[3][0] -= 2.0f * M_PI;
while (s_status[3][0] < -M_PI) s_status[3][0] += 2.0f * M_PI;
\end{lstlisting}
 
 
\subsubsection{Meetmodel en Kalman update}
 
De updatestap werkt de schatting bij op basis van het bewegingsmodel. Zodra een nieuwe GNSS-meting binnenkomt wordt de schatting gecorrigeerd op basis van het verschil tussen wat het filter verwachtte en wat de GNSS daadwerkelijk meet. Dit is de updatestap.
 
De GNSS-ontvanger levert elke 40 ms een positiemeting in het lokale ENU-stelsel. Niet alle zes toestandsgrootheden worden direct gemeten — alleen de positie $X$ en $Y$ zijn direct zichtbaar in de GNSS-data. Het meetmodel $H$ beschrijft welke toestandsgrootheden worden gemeten. Het is een $2 \times 6$ matrix die aangeeft dat de meting alleen de eerste twee elementen van de toestandsvector betreft:
 
\begin{equation}
    H = \begin{bmatrix}
        1 & 0 & 0 & 0 & 0 & 0 \\
        0 & 1 & 0 & 0 & 0 & 0
    \end{bmatrix}
    \label{eq:H}
\end{equation}
 
Het verschil tussen de gemeten positie en de voorspelde positie wordt de innovatie $\mathbf{y}$ genoemd. Dit is het signaal dat het filter gebruikt om zijn schatting bij te sturen:
 
\begin{equation}
    \mathbf{y} = \mathbf{z} - H \cdot \mathbf{x}^-
    \label{eq:innovatie}
\end{equation}
 
De innovatiecovariantie $S$ geeft aan hoe groot de verwachte spreiding van de innovatie is. Dit is de som van de onzekerheid in de voorspelling en de onzekerheid in de meting:
 
\begin{equation}
    S = H \cdot P^- \cdot H^T + R
    \label{eq:S}
\end{equation}
 
De Kalman gain $K$ bepaalt vervolgens hoeveel gewicht de meting krijgt ten opzichte van de updatestap. Als de GNSS-meting erg nauwkeurig is (kleine $R$) is $K$ groot en wordt de meting sterk gevolgd. Als de GNSS-onzekerheid groot is zoals bij Lk=0, is $K$ klein en vertrouwt het filter meer op zijn eigen bewegingsmodel:
 
\begin{equation}
    K = P^- \cdot H^T \cdot S^{-1}
    \label{eq:kalman_gain}
\end{equation}
 
Met de Kalman gain wordt de toestandsschatting gecorrigeerd:
 
\begin{equation}
    \mathbf{x}^+ = \mathbf{x}^- + K \cdot \mathbf{y}
    \label{eq:update}
\end{equation}
 
Tot slot wordt de covariantiematrix bijgewerkt om te reflecteren dat de onzekerheid na de update kleiner is. Dit gebeurt in de zogenaamde Joseph-vorm, die wiskundig equivalent is aan de standaard formulering maar numeriek stabieler is en garandeert dat $P$ positief semidefiniet blijft ook bij afrondingsfouten:
 
\begin{equation}
    P^+ = (I - K H) \cdot P^- \cdot (I - K H)^T + K \cdot R \cdot K^T
    \label{eq:joseph}
\end{equation}
 
    
\subsubsection{NIS-gebaseerde metingrejectie}
 
Een GNSS-ontvanger kan soms een sterk afwijkende positie rapporteren door bijvoorbeeld multipath of een tijdelijke False Fixed situatie. Als het filter zo'n meting klakkeloos verwerkt, kan de positieschatting in één stap meters verspringen. Om dit te voorkomen wordt elke meting vooraf getoetst op statistische plausibiliteit via de Normalized Innovation Squared.
 
De NIS meet hoe verrassend de binnenkomende meting is gegeven de huidige onzekerheid van het filter. Een grote NIS betekent dat de meting ver buiten de verwachte spreiding valt. De berekening is:
 
\begin{equation}
    \text{NIS} = \mathbf{y}^T \cdot S^{-1} \cdot \mathbf{y}
    \label{eq:NIS}
\end{equation}
 
Als het filter goed gekalibreerd is, volgt de NIS een chi-kwadraatverdeling met twee vrijheidsgraden — één voor elke gemeten dimensie. Bij een betrouwbaarheidsniveau van 99\% bedraagt de drempelwaarde $\chi^2_{2,\,0{,}99} = 9{,}210$. Elke meting met een hogere NIS wordt als statistisch onwaarschijnlijk beschouwd en verworpen, waarna de updatestap gewoon doorgaat zonder correctie:
 
\begin{verbatim}
if (nis > 9.210f) {
    return;
}
\end{verbatim}
 
Bij verwerping groeit $P$ via de volgende updatestap verder, waardoor de Kalman gain bij de eerstvolgende betrouwbare meting groter wordt en het filter sneller kan corrigeren. Naast zijn functie als uitschuiverdetectie is de NIS ook een diagnosemiddel voor de filterafstemming. Als de gemiddelde NIS structureel boven 2 ligt worden metingen vaker dan verwacht als onwaarschijnlijk bestempeld, wat duidt op te kleine ruismatrices. Als de gemiddelde NIS structureel onder 2 ligt vertrouwt het filter de metingen te sterk, wat duidt op te grote ruismatrices.
 
 
\subsubsection{Stilstanddetectie (Zero Velocity Update)}
 
Wanneer het systeem stilstaat maar de IMU blijft kleine versnellingsfouten meten, loopt de geschatte snelheid langzaam op terwijl die in werkelijkheid nul is. De Zero Velocity Update lost dit op door bij gedetecteerde stilstand een extra correctiestap uit te voeren die de snelheidsschatting terugtrekt naar nul. Dit werkt via dezelfde Kalman-mechanisme als de GNSS-update, maar in plaats van een positiemeting wordt een pseudo-meting van $v = 0$ gebruikt. De innovatie is simpelweg het verschil tussen de verwachte stilstand en de geschatte snelheid:
 
\begin{equation}
    \text{innovatie} = 0 - v_{\text{geschat}}
    \label{eq:zupt_innovatie}
\end{equation}
 
De bijbehorende Kalman gain wordt alleen berekend voor de snelheidsrij van $P$, zodat alleen de snelheid en de indirect gekoppelde grootheden worden bijgestuurd en de positieschatting onaangetast blijft:
 
\begin{verbatim}
float innov = 0.0f - s_status[2][0];
float S     = s_P[2][2] + zupt_R;
for (int i = 0; i < 6; i++)
    K[i] = s_P[i][2] / S;
\end{verbatim}
 
Een bijkomend voordeel is dat de biassen tijdens stilstand sneller convergeren. Omdat het filter bij stilstand zekerheid heeft over zowel de positie (via GNSS) als de snelheid (via ZUPT), kan het de resterende afwijkingen volledig toeschrijven aan de sensorbiassen en die daarmee nauwkeuriger bijschatten.
 
De stilstanddetectie is gebaseerd op een variantieanalyse over een schuivend venster van 32 IMU-samples, wat bij 200 Hz overeenkomt met een venster van 160 ms. Per venster worden de variantie en het gemiddelde berekend voor alle zes IMU-kanalen. Het systeem wordt pas als stilstaand beschouwd wanneer aan de volgende drie voorwaarden tegelijkertijd wordt voldaan:
 
\begin{enumerate}
    \item De variantie van de accelerometer is op alle drie assen kleiner dan de ingestelde drempelwaarde. Dit filtert kleine trillingen eruit die ook bij stilstand aanwezig zijn.
    \item De variantie van de gyroscoop is op alle drie assen kleiner dan de ingestelde drempelwaarde.
    \item De absolute waarde van het gemiddelde gyroscoopsignaal is op alle drie assen kleiner dan een absolute drempel. Deze voorwaarde voorkomt dat rechtdoor rijden met weinig variatie — waarbij de variantie laag is maar het systeem wel degelijk beweegt — als stilstand wordt geïnterpreteerd.
\end{enumerate}
 
Om valse activering te voorkomen zijn twee hysterese-timers geïmplementeerd. Stilstand wordt pas bevestigd na 1000 ms aaneengesloten stilstanddetectie, terwijl beweging al na 150 ms wordt herkend. Dit asymmetrische ontwerp zorgt ervoor dat een korte stilstaande periode zoals stoppen voor een obstakel niet direct als stilstand wordt geregistreerd, maar dat hervatting van de beweging snel wordt gedetecteerd.
\end{comment} 


\newpage